\section{A vanishing--variance rough obstruction}
\label{sec:tpc44-sharp-obstruction}

We now show that the high--prime condition left by the preceding
sections cannot be supplied by a black--box variance, fixed--moment, or
low--influence theorem.  The example belongs to the same abstract class
as \eqref{eq:tpc44-coordinate-packet}: it is the squared norm of a Walsh
sum with distinct squarefree labels.

Fix \(z\), choose distinct primes \(q_1,\ldots,q_R>z\), and choose a
squarefree \(u_0\) coprime to all of them.  For \(\lambda>0\), put
\begin{equation}
 S(\varepsilon)
 :=\chi_{u_0}(\varepsilon)
 \left(
  1-\frac{\lambda}{R}\sum_{i=1}^{R}\chi_{q_i}(\varepsilon)
 \right),
 \qquad Z(\varepsilon):=|S(\varepsilon)|^2.
 \label{eq:tpc44-obstruction-SZ}
\end{equation}
The distinct row labels are \(u_0,u_0q_1,\ldots,u_0q_R\).

\begin{theorem}[Rough degree--two late collapse]
\label{thm:tpc44-rough-obstruction}
The energy \eqref{eq:tpc44-obstruction-SZ} has the exact expansion
\begin{equation}
 Z
 =1+\frac{\lambda^2}{R}
 -\frac{2\lambda}{R}\sum_{i=1}^{R}\chi_{q_i}
 +\frac{2\lambda^2}{R^2}
  \sum_{1\le i<j\le R}\chi_{q_iq_j}.
 \label{eq:tpc44-obstruction-expansion}
\end{equation}
Thus every nonconstant kernel is \(z\)-rough and the Walsh degree is at
most \(2\).  Moreover,
\begin{align}
 D:=\E Z&=1+\frac{\lambda^2}{R},
 \label{eq:tpc44-obstruction-D}\\
 Z(-\one)&=(1+\lambda)^2,
 \label{eq:tpc44-obstruction-corner}\\
 \operatorname{Var}Z
 &=\frac{4\lambda^2}{R}
  +\frac{2\lambda^4(R-1)}{R^3}.
 \label{eq:tpc44-obstruction-variance}
\end{align}
Take \(\lambda=R^{1/4}\).  Then
\begin{equation}
 D\longrightarrow1,
 \qquad
 \frac{\operatorname{Var}Z}{D^2}=O(R^{-1/2})\longrightarrow0,
 \qquad
 \frac{Z(-\one)}D\sim R^{1/2}\longrightarrow\infty.
 \label{eq:tpc44-obstruction-asymptotics}
\end{equation}
For every fixed \(1\le q<\infty\),
\begin{equation}
 \left\|\frac ZD-1\right\|_{L^q(U)}\longrightarrow0.
 \label{eq:tpc44-obstruction-fixed-moments}
\end{equation}
If \(d\QZ=(Z/D)dU\), then
\begin{equation}
 \chi^2(\QZ\|U)=\frac{\operatorname{Var}Z}{D^2}\longrightarrow0,
 \qquad
 \|\QZ-U\|_{\rm TV}\longrightarrow0,
 \label{eq:tpc44-obstruction-TV}
\end{equation}
while the density at the distinguished point diverges as in
\eqref{eq:tpc44-obstruction-asymptotics}.
\end{theorem}

\begin{proof}
The factor \(\chi_{u_0}\) has modulus one and disappears on squaring.
Expanding the remaining real square, using
\(
 (\sum_i\chi_{q_i})^2=R+2\sum_{i<j}\chi_{q_iq_j}
\), gives \eqref{eq:tpc44-obstruction-expansion}.  Orthogonality gives
the mean and variance, while evaluating every \(\chi_{q_i}\) at \(-1\)
gives \eqref{eq:tpc44-obstruction-corner}.

For \(\lambda=R^{1/4}\), write
\begin{equation}
 \frac{\lambda}{R}\sum_{i=1}^{R}\chi_{q_i}
 =R^{-1/4}
  \left(R^{-1/2}\sum_{i=1}^{R}\chi_{q_i}\right).
 \label{eq:tpc44-obstruction-Khintchine-scaling}
\end{equation}
Khintchine's inequality bounds every fixed \(L^q\) norm of the
parenthesized variable.  Hence the left side tends to zero in every
fixed \(L^q\), and squaring proves
\eqref{eq:tpc44-obstruction-fixed-moments}.  The chi--square identity is
the definition of \(\QZ\), and total variation is at most one half of
its square root.
\end{proof}

Even the primewise face increments are individually negligible.  Order
the \(q_i\) arbitrarily and define \(A_t\) as in
\eqref{eq:tpc44-Ar}.  Direct averaging gives
\begin{equation}
 A_t=
 \left(1+\lambda\frac{R-t}{R}\right)^2
 +\frac{\lambda^2t}{R^2},
 \label{eq:tpc44-obstruction-At}
\end{equation}
and
\begin{equation}
 A_{t-1}-A_t
 =\frac{2\lambda}{R}
  \left(1+\lambda\frac{R-t}{R}\right).
 \label{eq:tpc44-obstruction-Lt}
\end{equation}
Consequently
\begin{equation}
 0\le\delta_t\le\frac{2\lambda}{R},
 \qquad
 \sum_{t=1}^{R}\delta_t^2\le\frac{4\lambda^2}{R}.
 \label{eq:tpc44-obstruction-small-deltas}
\end{equation}
For \(\lambda=R^{1/4}\), the maximum is \(O(R^{-3/4})\) and the sum of
squares is \(O(R^{-1/2})\), yet
\begin{equation}
 \prod_{t=1}^{R}(1+\delta_t)
 =\frac{(1+\lambda)^2}{1+\lambda^2/R}
 \sim R^{1/2}.
 \label{eq:tpc44-obstruction-delta-product}
\end{equation}

For completeness, the normalized Walsh influence of a prime \(q_i\)
is also \(O(R^{-1})\).  Indeed
\begin{equation}
 \operatorname{Inf}_{q_i}(F)
 :=\sum_{q_i\mid\kappa}|C(\kappa)|^2
 =\frac{4\lambda^2}{R^2}
  +\frac{4\lambda^4(R-1)}{R^4},
 \label{eq:tpc44-obstruction-influence}
\end{equation}
and division by \eqref{eq:tpc44-obstruction-variance} gives the claim at
\(\lambda=R^{1/4}\).

\begin{remark}[Scope of the counterexample]
\label{rem:tpc44-obstruction-scope}
The example proves that the abstract nonnegative Walsh--Gram interface
does not force good behavior at the all--minus corner.  It does not show
that the actual packet coefficient \(C_X(\kappa)\) realizes this example,
nor that the physical endpoint fails.  Its role is to identify which
additional arithmetic statement is indispensable.
\end{remark}
